\begin{tcolorbox}[colback=propbox, colframe=propborder, arc=2pt,
  left=6pt, right=6pt, top=4pt, bottom=4pt,
  title={\small\textbf{Definition (Atomic Formula)}},
  fonttitle=\small\bfseries]
\label{def:atomic-formula-propositional-logic}
An \emph{atomic formula} is a well-formed formula that is a propositional
variable.
\end{tcolorbox}

\begin{remark*}[Interpretation]
Atomic formulas are the base cases of the recursive grammar. They contain no
connective structure.
\end{remark*}

\begin{dependencies}
\begin{itemize}
  \item \hyperref[def:well-formed-formula]{Well-Formed Formula}
  \item \hyperref[def:propositional-variable]{Propositional Variable}
\end{itemize}
\end{dependencies}

\begin{tcolorbox}[colback=propbox, colframe=propborder, arc=2pt,
  left=6pt, right=6pt, top=4pt, bottom=4pt,
  title={\small\textbf{Definition (Molecular Formula)}},
  fonttitle=\small\bfseries]
\label{def:molecular-formula}
A \emph{molecular formula} is a well-formed formula that is not atomic;
equivalently, it is formed using at least one logical connective.
\end{tcolorbox}

\begin{remark*}[Interpretation]
Molecular formulas are exactly the formulas with nontrivial syntactic
construction. They arise from the negation clause or a binary connective clause.
\end{remark*}

\begin{dependencies}
\begin{itemize}
  \item \hyperref[def:well-formed-formula]{Well-Formed Formula}
  \item \hyperref[def:logical-connective]{Logical Connective}
\end{itemize}
\end{dependencies}
